% vim: set textwidth=78 autoindent:

\section{Uso dei Plugin esterni QGIS Python}\label{sec:external_plugins}\index{plugins}

% when the revision of a section has been finalized, 
% comment out the following line:
% \updatedisclaimer

I Plugin esterni QGIS sono scritti in python. Sono immagazzinati in un archivio ufficiale moderato e mantenuto dall'autore. La tabella \ref{tab:external_plugins} mostra una lista di plugin attualmente disponibile con una breve descrizione.\footnote{Anche gli aggiornamenti dei plugin core possono essere disponibili anche in questo archivio come overlay esterni.} 
\footnote{L'Istallatore di Plugin Python è anch'esso un plugin Python esterno, ma è parte della sorgente QGIS e caricato automaticamente e selezionabile nel Gestore dei Plugin di QGIS (vedere Sezione~\ref{sec:load_external_plugin}).}

Quando questo manuale è stato distribuito, archivio ufficiale moderato e mantenuto non era ancora completamente costituito. Una documentazione dettagliato circa l'uso, l'autore ed altre importanti informazioni sono fornite con il plugin esterno stesso e non è parte di questo manuale.

Si può trovare una lista aggirnata dei plugin esterni moderati nell'Archivio ufficiale di QGIS a \dropmenuopttwo{plugin_installer}{Scarica Python
Plugins...}) e \url{http://qgis.osgeo.org/download/plugins.html}. 

\begin{table}[H]
\centering
\caption{Attuali Plugin QGIS esterni moderati}\label{tab:external_plugins}\medskip
\small
 \begin{tabular}{|l|l|p{4in}|}
\hline \textbf{Icona} & \textbf{Plugin esterno} & \textbf{Descrizione}\\
\hline
\includegraphics[width=0.7cm]{zoom2point_icon}
 & Zoom al punto \index{plugins!Zoom To Point} & Fa lo zoom alle coordinate specificate nella finastra di dialogo di input. Si può specificare il livello di zoom così come il controllo dell'estensione della vista.\\
\hline
\end{tabular}
\end{table}

Una descrizione dettagliata per i plugin python esterni si può trovare nella 
Sezione \ref{sec:load_external_plugin}.

\minisec{Archivio Plugin Python sviluppati degli utenti e Archivi degli autori}

In aggiunta ai plugin esterni moderati, esiste anche un altro archivio Plugin Python non ufficiale. Contiene plugin che non sono ancora abbastanza maturi da essere iclusi nell'archivio ufficiale, comunque alcuni possono essere abbastanza utili. Inoltre, alcuni dei nostri contributori mantengono i propri archivi.

Per aggiungere l'archivio non ufficiale e gli archivi degli autori aprire l'Istallatore di Plugin Python (\mainmenuopt{Plugins} > \dropmenuopttwo{plugin_installer}{Scarico Plugin Python}),
andare alla scheda \tab{Repositories} e premere \button{Aggiunge reposotory di terze parti}. Se non si vuole uno o più degli archivi aggiunti, disabilitarli con il tasto
\button{Edit...} o rimuoverli completamente con il tasto \button{Delete}.

\begin{Tip} \caption{\textsc{Aggiungere più Plugin esterni}}
\qgistip{In aggiunta ai plugin esterni moderati dall'archivio ufficiale QGIS si possono aggiungere più archivi esterni. Quindi selezionare la scheda Repositories nell'Istallatore di Plugin Pythons}
\end{Tip}

